\FloatBarrier
\section{\Acl*{ITZ} (\acs*{ITZ})}

Cement based mortars and concrete are a mixture of aggregates and a cement matrix.
Around the interface of aggregates, walls or reinforcement and the binder a thin shell of a less dense matrix forms, which is called \ac{ITZ}.
It is caused by the \textit{wall effect}, which causes a depletion of larger cement grains in this area, resulting in a high-porosity zone \cite{gao.2014, torsaeter.2015}.
This higher porosity influences the mechanical properties \cite{zhu.2000,xiao.2013} and causes increased water permeability \cite{ye.2007} and is therefore the `weak link' in in cement composites \cite{zhu.2000}.
Therefore, this zone is prone to microcracks \cite{ye.2007}.

The \ac{ITZ} can be measured based on the analysis of the porosity around aggregates using \ac{SEM}-\ac{BSE} images \cite{gao.2014} or \ac{mCT} datasets \cite{lavrov.2017}, by measuring mechanical parameters like hardness or the \textsc{Young}'s modulus using \ac{NI} \cite{zhu.2000,xiao.2013} or by simulating the structure \cite{ye.2007}.
However, the width of the \ac{ITZ} varies significantly in the literature as the selection in \zcref{tab:ITZ-width} shows.

\begin{table}[htbp]
	\centering
		\caption{
			Overview of \acs*{ITZ} widths reported in literature depending on the method and interface type.
		}
		\label{tab:ITZ-width}
		\begin{tabularx}{\linewidth}{>{\raggedleft\arraybackslash}l>{\raggedright\arraybackslash}X>{\raggedright\arraybackslash}X>{\centering\arraybackslash}l}
			\toprule
			\acs*{ITZ} width & Interface type & Method & Reference \\
			\midrule
			\qtyrange{15}{20}{\micro\meter} & natural aggregates & calculated from \acs*{MIP} measurements & \cite{winslow.1994} \\
			\qtyrange{30}{100}{\micro\meter} & natural aggregates & \ac{SEM}-\ac{BSE} (pores) & \cite{scrivener.1996} \\
			\qtyrange{40}{50}{\micro\meter} & natural aggregates & \acs{NI} & \cite{xiao.2013} \\
			\qtyrange{30}{50}{\micro\meter} & natural aggregates & \ac{SEM}-\ac{BSE} (pores) & \cite{skarzynski.2016}\\
			$\approx$\,\qty{10}{\micro\meter} & natural aggregates & \ac{SEM}-\ac{BSE} (pores) & \cite{branch.2018} \\
			\qtyrange{38}{42}{\micro\meter} & natural aggregates & \ac{SEM}-\ac{EDX} (\ce{Ca}/\ce{Si} ratio) & \cite{tian.2018} \\
			\qty{40}{\micro\meter} & recycled aggregates & \acs{NI} & \cite{zhang.2019} \\
			<\,\qty{50}{\micro\meter} & natural aggregates & \ac{SEM}-\ac{BSE} (pores) & \cite{li.2026} \\
			\qtyrange{55}{65}{\micro\meter} & recycled aggregates & \acs{NI} & \cite{xiao.2013} \\
			\qtyrange{50}{60}{\micro\meter} & recycled aggregates & \acs{NI} & \cite{zhang.2019} \\
			<\,\qty{60}{\micro\meter} & recycled aggregates & \ac{SEM}-\ac{BSE} (pores) & \cite{li.2026} \\
			\qty{65}{\micro\meter} & recycled brick & \ac{SEM}-\ac{BSE} (pores) & \cite{li.2025} \\
			<\,\qty{100}{\micro\meter} & steel wall & \ac{SEM}-\ac{BSE} (grain) & \cite{torsaeter.2015} \\
			<\,\qty{30}{\micro\meter} & steel rebar & \acs{NI} & \cite{zhu.2000} \\
			<\,\qty{20}{\micro\meter} & wall & simulation (HYMOSTRUC3D) & \cite{ye.2007} \\
			<\,\qty{15}{\micro\meter} & wall & \ac{SEM}-\ac{BSE} (pores), simulation (HYMOSTRUC3D) & \cite{gao.2014} \\
			<\,\qty{20}{\micro\meter} & glass tube & \ac{mCT} & \cite{lavrov.2017} \\
			\bottomrule
		\end{tabularx}
\end{table}

The variation of the \ac{ITZ} width seems to be influenced by the interface material and the method used to measure it.
Using the porosity or the mechanical properties of the material are the most common.
Using the remaining unhydrated cement particle density as in \cite{torsaeter.2015} seems to overestimate the \ac{ITZ} width and is a unpractical method for real mortars or concretes.
A consistent observation in literature which compared recycled and natural aggregates was, that recycled aggregates cause a \qtyrange{10}{20}{\micro\meter} wider \ac{ITZ}.
Furthermore, combined \ac{EDX} analysis revealed a \ce{Ca}-rich layer directly on the interface of the aggregates \cite{li.2012, tian.2018}.
This points to the crystallisation of \ac{CH} in this area, which is supported by simulations using HYMOSTRUC3D \cite{gao.2014}.


